\documentclass[11pt,letterpaper]{article}
\usepackage[technical-reference]{cornell-notes}
\usepackage{needspace}

\title{Clause 20: Memory Management Library - Cornell Notes}
\author{}
\date{}
\setDocTitle{Clause 20: Memory Management Library}
\setDocSubtitle{C++ 2024 Cornell Notes}
\setDocOwner{C++ 2024 Cornell Notes}
\setDocAuthor{}
\setDocDate{}
\setCornellCollection{C++ 2024 Cornell Notes}
\setCornellUnitType{Clause}
\setCornellUnitNumber{20}
\setCornellUnitTitle{Memory Management Library}
\setCornellCompanionLabel{C++ Technical Reference}
\hypersetup{pdftitle={Clause 20: Memory Management Library - Cornell Notes},pdfsubject={C++ technical study notes},pdfauthor={}}

\begin{document}
\maketitle
\begin{CornellReferenceOrientationBox}
\textbf{Purpose.} Defines allocators, allocator traits, pointer traits, uninitialized memory algorithms, smart pointers, ownership utilities, polymorphic memory resources, and explicit lifetime helpers.
\par\textbf{Role.} The library model for acquiring storage, constructing objects, transferring ownership, and reclaiming resources.
\par\textbf{Principal dependencies.} Uses the library-wide rules of Clause 16 together with the applicable core-language concepts and supporting library clauses.
\par\textbf{Primary audience.} programmers, container authors, and allocator implementers.
\par\textbf{Major subject groups.} General; Memory; Smart pointers; Memory resources; Class template scoped\_\allowbreak{}allocator\_\allowbreak{}adaptor.
\end{CornellReferenceOrientationBox}
\section{Learning Objectives}
\begin{itemize}[label=\textcolor{CornellReferenceTeal}{$\blacktriangleright$}]
\item Define the governing ideas of General and apply its stated contract at a call site.
\item Identify the governing ideas of Memory and apply its stated contract at a call site.
\item Distinguish the governing ideas of Smart pointers and apply its stated contract at a call site.
\item Explain the governing ideas of Memory resources and apply its stated contract at a call site.
\item Trace the governing ideas of Class template scoped\_\allowbreak{}allocator\_\allowbreak{}adaptor and apply its stated contract at a call site.
\item Apply the governing ideas of General and apply its stated contract at a call site.
\item Analyze the governing ideas of Header <memory> synopsis and apply its stated contract at a call site.
\item Evaluate the governing ideas of Pointer traits and apply its stated contract at a call site.
\item Diagnose the governing ideas of Pointer conversion and apply its stated contract at a call site.
\item Compare the governing ideas of Pointer alignment and apply its stated contract at a call site.
\item Verify the governing ideas of Explicit lifetime management and apply its stated contract at a call site.
\item Relate the governing ideas of Allocator argument tag and apply its stated contract at a call site.
\item Classify the governing ideas of uses\_\allowbreak{}allocator and apply its stated contract at a call site.
\item Justify the governing ideas of Allocator traits and apply its stated contract at a call site.
\item Audit the governing ideas of The default allocator and apply its stated contract at a call site.
\item Summarize the governing ideas of addressof and apply its stated contract at a call site.
\end{itemize}
\section{Normative Structure Map}
\small\begin{longtable}{@{}>{\RaggedRight\arraybackslash}p{0.13\textwidth}>{\RaggedRight\arraybackslash}p{0.27\textwidth}>{\RaggedRight\arraybackslash}p{0.19\textwidth}>{\RaggedRight\arraybackslash}p{0.31\textwidth}@{}}
\toprule\textbf{Subclause} & \textbf{Subject} & \textbf{Rule category} & \textbf{Key dependencies / entities}\\\midrule\endfirsthead
\toprule\textbf{Subclause} & \textbf{Subject} & \textbf{Rule category} & \textbf{Key dependencies / entities}\\\midrule\endhead
\bottomrule\endfoot
20.1 & General & library semantics & Uses the library-wide rules of Clause 16 together with the applicable core-language concepts and supporting library clauses.\\
20.2 & Memory & library semantics & General, Header <memory> synopsis, Pointer traits, Pointer conversion, and related subclauses\\
20.3 & Smart pointers & library semantics & Unique-ownership pointers, Shared-ownership pointers, Smart pointer hash support, Smart pointer adaptors\\
20.4 & Memory resources & library semantics & Header <memory\_\allowbreak{}resource> synopsis, Class memory\_\allowbreak{}resource, Class template polymorphic\_\allowbreak{}allocator, Access to program-wide memory\_\allowbreak{}resource objects, and related subclauses\\
20.5 & Class template scoped\_\allowbreak{}allocator\_\allowbreak{}adaptor & library semantics & Header <scoped\_\allowbreak{}allocator> synopsis, Member types, Constructors, Members, and related subclauses\\
\end{longtable}\normalsize
\begin{CornellReferenceNormativeBox}A heading maps the clause; it does not replace the detailed conditions. For each operation, read requirements, effects, result, exceptions, complexity, remarks, and notes with their stated normative force.\end{CornellReferenceNormativeBox}
\subsection*{Rule Classification Guide}
\small\begin{longtable}{@{}>{\RaggedRight\arraybackslash}p{0.25\textwidth}>{\RaggedRight\arraybackslash}p{0.67\textwidth}@{}}\toprule\textbf{Label or category} & \textbf{How to read it}\\\midrule\endfirsthead\toprule\textbf{Label or category} & \textbf{How to read it (continued)}\\\midrule\endhead\bottomrule\endfoot
Constraints & Conditions controlling whether a declaration, overload, or specialization participates in the relevant context.\\
Mandates & Requirements checked for the described declaration or specialization; failure has the stated well-formedness consequence.\\
Preconditions & Obligations the caller establishes before evaluation; violating one forfeits the operation's contract.\\
Effects / Returns / Postconditions & Separate the state change, produced result, and state guaranteed after successful completion.\\
Throws / Error conditions & State how failure is reported and which exceptional paths are permitted or required.\\
Complexity & Bounds or exact counts for the operations named by the specification.\\
Remarks / Recommended practice & Remarks refine applicability or participation; recommended practice guides implementations without silently becoming a program requirement.\\
Exposition-only & A device for describing semantics, not a public name on which a program can depend.\\
\end{longtable}\normalsize
\section{Cornell Cue-and-Notes}
\Needspace{8\baselineskip}
\subsection*{20.1 General}
\begin{longtable}{@{}>{\RaggedRight\arraybackslash\columncolor{CornellReferenceCueGray}}p{0.28\textwidth}>{\RaggedRight\arraybackslash}p{0.64\textwidth}@{}}
\rowcolor{CornellReferenceNavy}\color{white}\textbf{Cue / recall prompt} & \color{white}\textbf{Notes}\\\endfirsthead
\rowcolor{CornellReferenceNavy}\color{white}\textbf{Cue / recall prompt} & \color{white}\textbf{Notes (continued)}\\\endhead
\arrayrulecolor{CornellReferenceLineGray}
\textbf{What contract governs General?}\\[-1mm]{\scriptsize\CornellClauseRef{20.1}} & Frames the terminology, applicability, and relationships needed to interpret General; detailed rules remain in the following subclauses.\\\hline
\end{longtable}
\begin{CornellReferenceSummaryBox}General supplies the library semantics portion of this clause. The key review move is to connect its local wording to the clause-wide model, then classify each condition by normative force before relying on it.\end{CornellReferenceSummaryBox}
\Needspace{8\baselineskip}
\subsection*{20.2 Memory}
\begin{longtable}{@{}>{\RaggedRight\arraybackslash\columncolor{CornellReferenceCueGray}}p{0.28\textwidth}>{\RaggedRight\arraybackslash}p{0.64\textwidth}@{}}
\rowcolor{CornellReferenceNavy}\color{white}\textbf{Cue / recall prompt} & \color{white}\textbf{Notes}\\\endfirsthead
\rowcolor{CornellReferenceNavy}\color{white}\textbf{Cue / recall prompt} & \color{white}\textbf{Notes (continued)}\\\endhead
\arrayrulecolor{CornellReferenceLineGray}
\textbf{What contract governs General?}\\[-1mm]{\scriptsize\CornellClauseRef{20.2.1}} & Frames the terminology, applicability, and relationships needed to interpret General; detailed rules remain in the following subclauses.\\\hline
\textbf{What contract governs Header <memory> synopsis?}\\[-1mm]{\scriptsize\CornellClauseRef{20.2.2}} & Catalogs the declarations made available by Header <memory> synopsis. The synopsis is an interface map; the detailed specifications supply constraints and behavior.\\\hline
\textbf{What contract governs Pointer traits?}\\[-1mm]{\scriptsize\CornellClauseRef{20.2.3}} & Specifies the facility family Pointer traits. Read its declarations together with mandates, preconditions, effects, results, exceptions, complexity, and lifetime rules.\\\hline
\textbf{When is Pointer conversion permitted?}\\[-1mm]{\scriptsize\CornellClauseRef{20.2.4}} & Specifies source and destination conditions for Pointer conversion, the resulting type or value category, and any information-loss, pointer, qualification, or lifetime consequence.\\\hline
\textbf{What contract governs Pointer alignment?}\\[-1mm]{\scriptsize\CornellClauseRef{20.2.5}} & Specifies the facility family Pointer alignment. Read its declarations together with mandates, preconditions, effects, results, exceptions, complexity, and lifetime rules.\\\hline
\textbf{When does Explicit lifetime management begin or end?}\\[-1mm]{\scriptsize\CornellClauseRef{20.2.6}} & Tracks storage and object usability for Explicit lifetime management; storage availability alone does not prove that an object's lifetime has begun or remains active.\\\hline
\textbf{What contract governs Allocator argument tag?}\\[-1mm]{\scriptsize\CornellClauseRef{20.2.7}} & Specifies the facility family Allocator argument tag. Read its declarations together with mandates, preconditions, effects, results, exceptions, complexity, and lifetime rules.\\\hline
\textbf{What contract governs uses\_\allowbreak{}allocator?}\\[-1mm]{\scriptsize\CornellClauseRef{20.2.8}} & Specifies the facility family uses\_\allowbreak{}allocator. Read its declarations together with mandates, preconditions, effects, results, exceptions, complexity, and lifetime rules.\\\hline
\textbf{What contract governs Allocator traits?}\\[-1mm]{\scriptsize\CornellClauseRef{20.2.9}} & Specifies the facility family Allocator traits. Read its declarations together with mandates, preconditions, effects, results, exceptions, complexity, and lifetime rules.\\\hline
\textbf{What contract governs The default allocator?}\\[-1mm]{\scriptsize\CornellClauseRef{20.2.10}} & Specifies the facility family The default allocator. Read its declarations together with mandates, preconditions, effects, results, exceptions, complexity, and lifetime rules.\\\hline
\textbf{What contract governs addressof?}\\[-1mm]{\scriptsize\CornellClauseRef{20.2.11}} & Specifies the facility family addressof. Read its declarations together with mandates, preconditions, effects, results, exceptions, complexity, and lifetime rules.\\\hline
\textbf{What contract governs C library memory allocation?}\\[-1mm]{\scriptsize\CornellClauseRef{20.2.12}} & Specifies the facility family C library memory allocation. Read its declarations together with mandates, preconditions, effects, results, exceptions, complexity, and lifetime rules.\\\hline
\end{longtable}
\begin{CornellReferenceSummaryBox}Memory supplies the library semantics portion of this clause. The key review move is to connect its local wording to the clause-wide model, then classify each condition by normative force before relying on it.\end{CornellReferenceSummaryBox}
\Needspace{5\baselineskip}\paragraph{Detailed coverage ledger.}
\begin{description}[style=nextline,leftmargin=2.8cm,font=\normalfont\bfseries\color{CornellReferenceTeal}]
\item[20.2.3.1] \textbf{General.} Frames the terminology, applicability, and relationships needed to interpret General; detailed rules remain in the following subclauses.
\item[20.2.3.2] \textbf{Member types.} Specifies the facility family Member types. Read its declarations together with mandates, preconditions, effects, results, exceptions, complexity, and lifetime rules.
\item[20.2.3.3] \textbf{Member functions.} Specifies the facility family Member functions. Read its declarations together with mandates, preconditions, effects, results, exceptions, complexity, and lifetime rules.
\item[20.2.3.4] \textbf{Optional members.} Specifies the facility family Optional members. Read its declarations together with mandates, preconditions, effects, results, exceptions, complexity, and lifetime rules.
\item[20.2.8.1] \textbf{uses\_\allowbreak{}allocator trait.} Specifies the facility family uses\_\allowbreak{}allocator trait. Read its declarations together with mandates, preconditions, effects, results, exceptions, complexity, and lifetime rules.
\item[20.2.8.2] \textbf{Uses-allocator construction.} Specifies the facility family Uses-allocator construction. Read its declarations together with mandates, preconditions, effects, results, exceptions, complexity, and lifetime rules.
\item[20.2.9.1] \textbf{General.} Frames the terminology, applicability, and relationships needed to interpret General; detailed rules remain in the following subclauses.
\item[20.2.9.2] \textbf{Member types.} Specifies the facility family Member types. Read its declarations together with mandates, preconditions, effects, results, exceptions, complexity, and lifetime rules.
\item[20.2.9.3] \textbf{Static member functions.} Specifies the facility family Static member functions. Read its declarations together with mandates, preconditions, effects, results, exceptions, complexity, and lifetime rules.
\item[20.2.9.4] \textbf{Other.} Specifies the facility family Other. Read its declarations together with mandates, preconditions, effects, results, exceptions, complexity, and lifetime rules.
\item[20.2.10.1] \textbf{General.} Frames the terminology, applicability, and relationships needed to interpret General; detailed rules remain in the following subclauses.
\item[20.2.10.2] \textbf{Members.} Specifies the facility family Members. Read its declarations together with mandates, preconditions, effects, results, exceptions, complexity, and lifetime rules.
\item[20.2.10.3] \textbf{Operators.} Specifies the facility family Operators. Read its declarations together with mandates, preconditions, effects, results, exceptions, complexity, and lifetime rules.
\end{description}
\Needspace{8\baselineskip}
\subsection*{20.3 Smart pointers}
\begin{longtable}{@{}>{\RaggedRight\arraybackslash\columncolor{CornellReferenceCueGray}}p{0.28\textwidth}>{\RaggedRight\arraybackslash}p{0.64\textwidth}@{}}
\rowcolor{CornellReferenceNavy}\color{white}\textbf{Cue / recall prompt} & \color{white}\textbf{Notes}\\\endfirsthead
\rowcolor{CornellReferenceNavy}\color{white}\textbf{Cue / recall prompt} & \color{white}\textbf{Notes (continued)}\\\endhead
\arrayrulecolor{CornellReferenceLineGray}
\textbf{What contract governs Unique-ownership pointers?}\\[-1mm]{\scriptsize\CornellClauseRef{20.3.1}} & Specifies the facility family Unique-ownership pointers. Read its declarations together with mandates, preconditions, effects, results, exceptions, complexity, and lifetime rules.\\\hline
\textbf{What contract governs Shared-ownership pointers?}\\[-1mm]{\scriptsize\CornellClauseRef{20.3.2}} & Specifies the facility family Shared-ownership pointers. Read its declarations together with mandates, preconditions, effects, results, exceptions, complexity, and lifetime rules.\\\hline
\textbf{What contract governs Smart pointer hash support?}\\[-1mm]{\scriptsize\CornellClauseRef{20.3.3}} & Specifies the facility family Smart pointer hash support. Read its declarations together with mandates, preconditions, effects, results, exceptions, complexity, and lifetime rules.\\\hline
\textbf{What contract governs Smart pointer adaptors?}\\[-1mm]{\scriptsize\CornellClauseRef{20.3.4}} & Specifies the facility family Smart pointer adaptors. Read its declarations together with mandates, preconditions, effects, results, exceptions, complexity, and lifetime rules.\\\hline
\end{longtable}
\begin{CornellReferenceSummaryBox}Smart pointers supplies the library semantics portion of this clause. The key review move is to connect its local wording to the clause-wide model, then classify each condition by normative force before relying on it.\end{CornellReferenceSummaryBox}
\Needspace{5\baselineskip}\paragraph{Detailed coverage ledger.}
\begin{description}[style=nextline,leftmargin=2.8cm,font=\normalfont\bfseries\color{CornellReferenceTeal}]
\item[20.3.1.1] \textbf{General.} Frames the terminology, applicability, and relationships needed to interpret General; detailed rules remain in the following subclauses.
\item[20.3.1.2] \textbf{Default deleters.} Specifies the facility family Default deleters. Read its declarations together with mandates, preconditions, effects, results, exceptions, complexity, and lifetime rules.
\item[20.3.1.2.1] \textbf{General.} Frames the terminology, applicability, and relationships needed to interpret General; detailed rules remain in the following subclauses.
\item[20.3.1.2.2] \textbf{default\_\allowbreak{}delete.} Specifies the facility family default\_\allowbreak{}delete. Read its declarations together with mandates, preconditions, effects, results, exceptions, complexity, and lifetime rules.
\item[20.3.1.2.3] \textbf{default\_\allowbreak{}delete<T[]>.} Specifies the facility family default\_\allowbreak{}delete<T[]>. Read its declarations together with mandates, preconditions, effects, results, exceptions, complexity, and lifetime rules.
\item[20.3.1.3] \textbf{unique\_\allowbreak{}ptr for single objects.} Specifies the facility family unique\_\allowbreak{}ptr for single objects. Read its declarations together with mandates, preconditions, effects, results, exceptions, complexity, and lifetime rules.
\item[20.3.1.3.1] \textbf{General.} Frames the terminology, applicability, and relationships needed to interpret General; detailed rules remain in the following subclauses.
\item[20.3.1.3.2] \textbf{Constructors.} Specifies how Constructors establishes object state, including participation constraints, initialization effects, and exceptional exits.
\item[20.3.1.3.3] \textbf{Destructor.} Specifies how Destructor ends the object's managed state and which cleanup or termination consequences apply.
\item[20.3.1.3.4] \textbf{Assignment.} Governs state transfer or exchange for Assignment; check self-operation, validity after movement, exception behavior, and invalidation.
\item[20.3.1.3.5] \textbf{Observers.} Specifies the facility family Observers. Read its declarations together with mandates, preconditions, effects, results, exceptions, complexity, and lifetime rules.
\item[20.3.1.3.6] \textbf{Modifiers.} Specifies the facility family Modifiers. Read its declarations together with mandates, preconditions, effects, results, exceptions, complexity, and lifetime rules.
\item[20.3.1.4] \textbf{unique\_\allowbreak{}ptr for array objects with a runtime length.} Specifies the facility family unique\_\allowbreak{}ptr for array objects with a runtime length. Read its declarations together with mandates, preconditions, effects, results, exceptions, complexity, and lifetime rules.
\item[20.3.1.4.1] \textbf{General.} Frames the terminology, applicability, and relationships needed to interpret General; detailed rules remain in the following subclauses.
\item[20.3.1.4.2] \textbf{Constructors.} Specifies how Constructors establishes object state, including participation constraints, initialization effects, and exceptional exits.
\item[20.3.1.4.3] \textbf{Assignment.} Governs state transfer or exchange for Assignment; check self-operation, validity after movement, exception behavior, and invalidation.
\item[20.3.1.4.4] \textbf{Observers.} Specifies the facility family Observers. Read its declarations together with mandates, preconditions, effects, results, exceptions, complexity, and lifetime rules.
\item[20.3.1.4.5] \textbf{Modifiers.} Specifies the facility family Modifiers. Read its declarations together with mandates, preconditions, effects, results, exceptions, complexity, and lifetime rules.
\item[20.3.1.5] \textbf{Creation.} Specifies the facility family Creation. Read its declarations together with mandates, preconditions, effects, results, exceptions, complexity, and lifetime rules.
\item[20.3.1.6] \textbf{Specialized algorithms.} Defines the observable result of Specialized algorithms together with applicable iterator-domain, ordering, complexity, invalidation, and exception conditions.
\item[20.3.1.7] \textbf{I/O.} Specifies the facility family I/O. Read its declarations together with mandates, preconditions, effects, results, exceptions, complexity, and lifetime rules.
\item[20.3.2.1] \textbf{Class bad\_\allowbreak{}weak\_\allowbreak{}ptr.} Specifies the facility family Class bad\_\allowbreak{}weak\_\allowbreak{}ptr. Read its declarations together with mandates, preconditions, effects, results, exceptions, complexity, and lifetime rules.
\item[20.3.2.2] \textbf{Class template shared\_\allowbreak{}ptr.} Specifies the facility family Class template shared\_\allowbreak{}ptr. Read its declarations together with mandates, preconditions, effects, results, exceptions, complexity, and lifetime rules.
\item[20.3.2.2.1] \textbf{General.} Frames the terminology, applicability, and relationships needed to interpret General; detailed rules remain in the following subclauses.
\item[20.3.2.2.2] \textbf{Constructors.} Specifies how Constructors establishes object state, including participation constraints, initialization effects, and exceptional exits.
\item[20.3.2.2.3] \textbf{Destructor.} Specifies how Destructor ends the object's managed state and which cleanup or termination consequences apply.
\item[20.3.2.2.4] \textbf{Assignment.} Governs state transfer or exchange for Assignment; check self-operation, validity after movement, exception behavior, and invalidation.
\item[20.3.2.2.5] \textbf{Modifiers.} Specifies the facility family Modifiers. Read its declarations together with mandates, preconditions, effects, results, exceptions, complexity, and lifetime rules.
\item[20.3.2.2.6] \textbf{Observers.} Specifies the facility family Observers. Read its declarations together with mandates, preconditions, effects, results, exceptions, complexity, and lifetime rules.
\item[20.3.2.2.7] \textbf{Creation.} Specifies the facility family Creation. Read its declarations together with mandates, preconditions, effects, results, exceptions, complexity, and lifetime rules.
\item[20.3.2.2.8] \textbf{Comparison.} Specifies the facility family Comparison. Read its declarations together with mandates, preconditions, effects, results, exceptions, complexity, and lifetime rules.
\item[20.3.2.2.9] \textbf{Specialized algorithms.} Defines the observable result of Specialized algorithms together with applicable iterator-domain, ordering, complexity, invalidation, and exception conditions.
\item[20.3.2.2.10] \textbf{Casts.} Specifies the facility family Casts. Read its declarations together with mandates, preconditions, effects, results, exceptions, complexity, and lifetime rules.
\item[20.3.2.2.11] \textbf{get\_\allowbreak{}deleter.} Specifies the facility family get\_\allowbreak{}deleter. Read its declarations together with mandates, preconditions, effects, results, exceptions, complexity, and lifetime rules.
\item[20.3.2.2.12] \textbf{I/O.} Specifies the facility family I/O. Read its declarations together with mandates, preconditions, effects, results, exceptions, complexity, and lifetime rules.
\item[20.3.2.3] \textbf{Class template weak\_\allowbreak{}ptr.} Specifies the facility family Class template weak\_\allowbreak{}ptr. Read its declarations together with mandates, preconditions, effects, results, exceptions, complexity, and lifetime rules.
\item[20.3.2.3.1] \textbf{General.} Frames the terminology, applicability, and relationships needed to interpret General; detailed rules remain in the following subclauses.
\item[20.3.2.3.2] \textbf{Constructors.} Specifies how Constructors establishes object state, including participation constraints, initialization effects, and exceptional exits.
\item[20.3.2.3.3] \textbf{Destructor.} Specifies how Destructor ends the object's managed state and which cleanup or termination consequences apply.
\item[20.3.2.3.4] \textbf{Assignment.} Governs state transfer or exchange for Assignment; check self-operation, validity after movement, exception behavior, and invalidation.
\item[20.3.2.3.5] \textbf{Modifiers.} Specifies the facility family Modifiers. Read its declarations together with mandates, preconditions, effects, results, exceptions, complexity, and lifetime rules.
\item[20.3.2.3.6] \textbf{Observers.} Specifies the facility family Observers. Read its declarations together with mandates, preconditions, effects, results, exceptions, complexity, and lifetime rules.
\item[20.3.2.3.7] \textbf{Specialized algorithms.} Defines the observable result of Specialized algorithms together with applicable iterator-domain, ordering, complexity, invalidation, and exception conditions.
\item[20.3.2.4] \textbf{Class template owner\_\allowbreak{}less.} Specifies the facility family Class template owner\_\allowbreak{}less. Read its declarations together with mandates, preconditions, effects, results, exceptions, complexity, and lifetime rules.
\item[20.3.2.5] \textbf{Class template enable\_\allowbreak{}shared\_\allowbreak{}from\_\allowbreak{}this.} Specifies the facility family Class template enable\_\allowbreak{}shared\_\allowbreak{}from\_\allowbreak{}this. Read its declarations together with mandates, preconditions, effects, results, exceptions, complexity, and lifetime rules.
\item[20.3.4.1] \textbf{Class template out\_\allowbreak{}ptr\_\allowbreak{}t.} Specifies the facility family Class template out\_\allowbreak{}ptr\_\allowbreak{}t. Read its declarations together with mandates, preconditions, effects, results, exceptions, complexity, and lifetime rules.
\item[20.3.4.2] \textbf{Function template out\_\allowbreak{}ptr.} Specifies the facility family Function template out\_\allowbreak{}ptr. Read its declarations together with mandates, preconditions, effects, results, exceptions, complexity, and lifetime rules.
\item[20.3.4.3] \textbf{Class template inout\_\allowbreak{}ptr\_\allowbreak{}t.} Specifies the facility family Class template inout\_\allowbreak{}ptr\_\allowbreak{}t. Read its declarations together with mandates, preconditions, effects, results, exceptions, complexity, and lifetime rules.
\item[20.3.4.4] \textbf{Function template inout\_\allowbreak{}ptr.} Specifies the facility family Function template inout\_\allowbreak{}ptr. Read its declarations together with mandates, preconditions, effects, results, exceptions, complexity, and lifetime rules.
\end{description}
\Needspace{8\baselineskip}
\subsection*{20.4 Memory resources}
\begin{longtable}{@{}>{\RaggedRight\arraybackslash\columncolor{CornellReferenceCueGray}}p{0.28\textwidth}>{\RaggedRight\arraybackslash}p{0.64\textwidth}@{}}
\rowcolor{CornellReferenceNavy}\color{white}\textbf{Cue / recall prompt} & \color{white}\textbf{Notes}\\\endfirsthead
\rowcolor{CornellReferenceNavy}\color{white}\textbf{Cue / recall prompt} & \color{white}\textbf{Notes (continued)}\\\endhead
\arrayrulecolor{CornellReferenceLineGray}
\textbf{What contract governs Header <memory\_\allowbreak{}resource> synopsis?}\\[-1mm]{\scriptsize\CornellClauseRef{20.4.1}} & Catalogs the declarations made available by Header <memory\_\allowbreak{}resource> synopsis. The synopsis is an interface map; the detailed specifications supply constraints and behavior.\\\hline
\textbf{What contract governs Class memory\_\allowbreak{}resource?}\\[-1mm]{\scriptsize\CornellClauseRef{20.4.2}} & Specifies the facility family Class memory\_\allowbreak{}resource. Read its declarations together with mandates, preconditions, effects, results, exceptions, complexity, and lifetime rules.\\\hline
\textbf{What contract governs Class template polymorphic\_\allowbreak{}allocator?}\\[-1mm]{\scriptsize\CornellClauseRef{20.4.3}} & Specifies the facility family Class template polymorphic\_\allowbreak{}allocator. Read its declarations together with mandates, preconditions, effects, results, exceptions, complexity, and lifetime rules.\\\hline
\textbf{What contract governs Access to program-wide memory\_\allowbreak{}resource objects?}\\[-1mm]{\scriptsize\CornellClauseRef{20.4.4}} & Specifies the facility family Access to program-wide memory\_\allowbreak{}resource objects. Read its declarations together with mandates, preconditions, effects, results, exceptions, complexity, and lifetime rules.\\\hline
\textbf{What contract governs Pool resource classes?}\\[-1mm]{\scriptsize\CornellClauseRef{20.4.5}} & Specifies the facility family Pool resource classes. Read its declarations together with mandates, preconditions, effects, results, exceptions, complexity, and lifetime rules.\\\hline
\textbf{What contract governs Class monotonic\_\allowbreak{}buffer\_\allowbreak{}resource?}\\[-1mm]{\scriptsize\CornellClauseRef{20.4.6}} & Specifies the facility family Class monotonic\_\allowbreak{}buffer\_\allowbreak{}resource. Read its declarations together with mandates, preconditions, effects, results, exceptions, complexity, and lifetime rules.\\\hline
\end{longtable}
\begin{CornellReferenceSummaryBox}Memory resources supplies the library semantics portion of this clause. The key review move is to connect its local wording to the clause-wide model, then classify each condition by normative force before relying on it.\end{CornellReferenceSummaryBox}
\Needspace{5\baselineskip}\paragraph{Detailed coverage ledger.}
\begin{description}[style=nextline,leftmargin=2.8cm,font=\normalfont\bfseries\color{CornellReferenceTeal}]
\item[20.4.2.1] \textbf{General.} Frames the terminology, applicability, and relationships needed to interpret General; detailed rules remain in the following subclauses.
\item[20.4.2.2] \textbf{Public member functions.} Specifies the facility family Public member functions. Read its declarations together with mandates, preconditions, effects, results, exceptions, complexity, and lifetime rules.
\item[20.4.2.3] \textbf{Private virtual member functions.} Specifies the facility family Private virtual member functions. Read its declarations together with mandates, preconditions, effects, results, exceptions, complexity, and lifetime rules.
\item[20.4.2.4] \textbf{Equality.} Specifies the facility family Equality. Read its declarations together with mandates, preconditions, effects, results, exceptions, complexity, and lifetime rules.
\item[20.4.3.1] \textbf{General.} Frames the terminology, applicability, and relationships needed to interpret General; detailed rules remain in the following subclauses.
\item[20.4.3.2] \textbf{Constructors.} Specifies how Constructors establishes object state, including participation constraints, initialization effects, and exceptional exits.
\item[20.4.3.3] \textbf{Member functions.} Specifies the facility family Member functions. Read its declarations together with mandates, preconditions, effects, results, exceptions, complexity, and lifetime rules.
\item[20.4.3.4] \textbf{Equality.} Specifies the facility family Equality. Read its declarations together with mandates, preconditions, effects, results, exceptions, complexity, and lifetime rules.
\item[20.4.5.1] \textbf{Classes synchronized\_\allowbreak{}pool\_\allowbreak{}resource and unsynchronized\_\allowbreak{}pool\_\allowbreak{}resource.} Specifies the facility family Classes synchronized\_\allowbreak{}pool\_\allowbreak{}resource and unsynchronized\_\allowbreak{}pool\_\allowbreak{}resource. Read its declarations together with mandates, preconditions, effects, results, exceptions, complexity, and lifetime rules.
\item[20.4.5.2] \textbf{pool\_\allowbreak{}options data members.} Specifies the facility family pool\_\allowbreak{}options data members. Read its declarations together with mandates, preconditions, effects, results, exceptions, complexity, and lifetime rules.
\item[20.4.5.3] \textbf{Constructors and destructors.} Specifies how Constructors and destructors establishes object state, including participation constraints, initialization effects, and exceptional exits.
\item[20.4.5.4] \textbf{Members.} Specifies the facility family Members. Read its declarations together with mandates, preconditions, effects, results, exceptions, complexity, and lifetime rules.
\item[20.4.6.1] \textbf{General.} Frames the terminology, applicability, and relationships needed to interpret General; detailed rules remain in the following subclauses.
\item[20.4.6.2] \textbf{Constructors and destructor.} Specifies how Constructors and destructor establishes object state, including participation constraints, initialization effects, and exceptional exits.
\item[20.4.6.3] \textbf{Members.} Specifies the facility family Members. Read its declarations together with mandates, preconditions, effects, results, exceptions, complexity, and lifetime rules.
\end{description}
\Needspace{8\baselineskip}
\subsection*{20.5 Class template scoped\_\allowbreak{}allocator\_\allowbreak{}adaptor}
\begin{longtable}{@{}>{\RaggedRight\arraybackslash\columncolor{CornellReferenceCueGray}}p{0.28\textwidth}>{\RaggedRight\arraybackslash}p{0.64\textwidth}@{}}
\rowcolor{CornellReferenceNavy}\color{white}\textbf{Cue / recall prompt} & \color{white}\textbf{Notes}\\\endfirsthead
\rowcolor{CornellReferenceNavy}\color{white}\textbf{Cue / recall prompt} & \color{white}\textbf{Notes (continued)}\\\endhead
\arrayrulecolor{CornellReferenceLineGray}
\textbf{What contract governs Header <scoped\_\allowbreak{}allocator> synopsis?}\\[-1mm]{\scriptsize\CornellClauseRef{20.5.1}} & Catalogs the declarations made available by Header <scoped\_\allowbreak{}allocator> synopsis. The synopsis is an interface map; the detailed specifications supply constraints and behavior.\\\hline
\textbf{What contract governs Member types?}\\[-1mm]{\scriptsize\CornellClauseRef{20.5.2}} & Specifies the facility family Member types. Read its declarations together with mandates, preconditions, effects, results, exceptions, complexity, and lifetime rules.\\\hline
\textbf{How is Constructors initialized?}\\[-1mm]{\scriptsize\CornellClauseRef{20.5.3}} & Specifies how Constructors establishes object state, including participation constraints, initialization effects, and exceptional exits.\\\hline
\textbf{What contract governs Members?}\\[-1mm]{\scriptsize\CornellClauseRef{20.5.4}} & Specifies the facility family Members. Read its declarations together with mandates, preconditions, effects, results, exceptions, complexity, and lifetime rules.\\\hline
\textbf{What contract governs Operators?}\\[-1mm]{\scriptsize\CornellClauseRef{20.5.5}} & Specifies the facility family Operators. Read its declarations together with mandates, preconditions, effects, results, exceptions, complexity, and lifetime rules.\\\hline
\end{longtable}
\begin{CornellReferenceSummaryBox}Class template scoped\_\allowbreak{}allocator\_\allowbreak{}adaptor supplies the library semantics portion of this clause. The key review move is to connect its local wording to the clause-wide model, then classify each condition by normative force before relying on it.\end{CornellReferenceSummaryBox}
\section{Language-Lawyer and Library-Usage Pitfalls}
\begin{CornellReferencePitfallBox}\begin{itemize}[label=\textcolor{CornellReferenceAmber}{$\blacktriangleright$}]
\item Reading a synopsis as the complete behavioral contract.
\item Ignoring mandates, preconditions, exception behavior, or complexity requirements.
\item Using an exposition-only name as though it were public API.
\item Assuming invalidation, ownership, or thread-safety properties that are not stated.
\item Treating successful compilation as proof that semantic requirements and preconditions are met.
\end{itemize}\end{CornellReferencePitfallBox}
\section{Programmer and Implementation Checklist}
\begin{CornellReferenceChecklistBox}\begin{itemize}[label=$\square$]
\item Identify the declaring header or module exposure for each facility.
\item Check template arguments and mandates before evaluating an operation.
\item Establish every precondition at the call site.
\item Record effects, returned value, postconditions, and exceptions separately.
\item Audit iterator, reference, pointer, and object lifetime after mutation.
\item Verify stated complexity and synchronization assumptions.
\item For \CornellClauseRef{20.1}, verify general against the precise rule category and all cited dependencies.
\item For \CornellClauseRef{20.2}, verify memory against the precise rule category and all cited dependencies.
\item For \CornellClauseRef{20.3}, verify smart pointers against the precise rule category and all cited dependencies.
\item For \CornellClauseRef{20.4}, verify memory resources against the precise rule category and all cited dependencies.
\end{itemize}\end{CornellReferenceChecklistBox}
\section{Clause Synthesis}
\begin{CornellReferenceOrientationBox}Defines allocators, allocator traits, pointer traits, uninitialized memory algorithms, smart pointers, ownership utilities, polymorphic memory resources, and explicit lifetime helpers. The library model for acquiring storage, constructing objects, transferring ownership, and reclaiming resources. A sound review distinguishes syntax and participation from runtime preconditions, keeps notes and examples non-mandatory, and follows cross-clause dependencies before making a portability claim.\end{CornellReferenceOrientationBox}
\section{Self-Test Questions}
\begin{enumerate}
\item What role does General play in the clause's overall model?
\item Which normative distinctions must be kept separate when applying Memory?
\item How would you audit a program or implementation that depends on Smart pointers?
\item Compare Memory resources with Class template scoped\_\allowbreak{}allocator\_\allowbreak{}adaptor; what different question does each answer?
\item What role does Class template scoped\_\allowbreak{}allocator\_\allowbreak{}adaptor play in the clause's overall model?
\item Which normative distinctions must be kept separate when applying General?
\item How would you audit a program or implementation that depends on Header <memory> synopsis?
\item Compare Pointer traits with Pointer conversion; what different question does each answer?
\item What role does Pointer conversion play in the clause's overall model?
\item Which normative distinctions must be kept separate when applying Pointer alignment?
\item How would you audit a program or implementation that depends on Explicit lifetime management?
\item Compare Allocator argument tag with uses\_\allowbreak{}allocator; what different question does each answer?
\item What role does uses\_\allowbreak{}allocator play in the clause's overall model?
\item Which normative distinctions must be kept separate when applying Allocator traits?
\item Scenario: a program relies on an unstated implementation behavior while using The default allocator. What must be checked before calling the program portable?
\item Scenario: a program relies on an unstated implementation behavior while using addressof. What must be checked before calling the program portable?
\end{enumerate}
\clearpage\section{Answer Key}
\begin{enumerate}
\item General is one part of the clause's library semantics model. Its role is understood by connecting the local rule in 20.4.6.1 to the clause orientation and its dependent concepts.
\item Keep formation and well-formedness separate from runtime preconditions and effects. Also distinguish mandatory wording from notes, implementation choice from undefined behavior, and interface declarations from detailed contracts; see 20.2.
\item Audit Smart pointers by locating the controlling wording in 20.3, listing inputs and type constraints, proving any preconditions, recording effects and failure behavior, and checking lifetime, invalidation, complexity, and synchronization where applicable.
\item Memory resources and Class template scoped\_\allowbreak{}allocator\_\allowbreak{}adaptor occupy different steps or facility groups. Analyze each under its own subclause, then follow the explicit dependency between them rather than transferring one set of guarantees to the other; begin at 20.4.
\item Class template scoped\_\allowbreak{}allocator\_\allowbreak{}adaptor is one part of the clause's library semantics model. Its role is understood by connecting the local rule in 20.5 to the clause orientation and its dependent concepts.
\item Keep formation and well-formedness separate from runtime preconditions and effects. Also distinguish mandatory wording from notes, implementation choice from undefined behavior, and interface declarations from detailed contracts; see 20.4.6.1.
\item Audit Header <memory> synopsis by locating the controlling wording in 20.2.2, listing inputs and type constraints, proving any preconditions, recording effects and failure behavior, and checking lifetime, invalidation, complexity, and synchronization where applicable.
\item Pointer traits and Pointer conversion occupy different steps or facility groups. Analyze each under its own subclause, then follow the explicit dependency between them rather than transferring one set of guarantees to the other; begin at 20.2.3.
\item Pointer conversion is one part of the clause's library semantics model. Its role is understood by connecting the local rule in 20.2.4 to the clause orientation and its dependent concepts.
\item Keep formation and well-formedness separate from runtime preconditions and effects. Also distinguish mandatory wording from notes, implementation choice from undefined behavior, and interface declarations from detailed contracts; see 20.2.5.
\item Audit Explicit lifetime management by locating the controlling wording in 20.2.6, listing inputs and type constraints, proving any preconditions, recording effects and failure behavior, and checking lifetime, invalidation, complexity, and synchronization where applicable.
\item Allocator argument tag and uses\_\allowbreak{}allocator occupy different steps or facility groups. Analyze each under its own subclause, then follow the explicit dependency between them rather than transferring one set of guarantees to the other; begin at 20.2.7.
\item uses\_\allowbreak{}allocator is one part of the clause's library semantics model. Its role is understood by connecting the local rule in 20.2.8 to the clause orientation and its dependent concepts.
\item Keep formation and well-formedness separate from runtime preconditions and effects. Also distinguish mandatory wording from notes, implementation choice from undefined behavior, and interface declarations from detailed contracts; see 20.2.9.
\item First identify the exact requirement in 20.2.10, then classify the relied-on behavior as required, unspecified, implementation-defined, conditionally supported, or undefined. Portability requires satisfying the program obligations and avoiding dependence on a choice not guaranteed in the target environments.
\item First identify the exact requirement in 20.2.11, then classify the relied-on behavior as required, unspecified, implementation-defined, conditionally supported, or undefined. Portability requires satisfying the program obligations and avoiding dependence on a choice not guaranteed in the target environments.
\end{enumerate}
\section{Key-Term Recap}
\begin{longtable}{@{}>{\RaggedRight\arraybackslash}p{0.28\textwidth}>{\RaggedRight\arraybackslash}p{0.64\textwidth}@{}}\toprule\textbf{Term} & \textbf{Working meaning}\\\midrule\endfirsthead\toprule\textbf{Term} & \textbf{Working meaning (continued)}\\\midrule\endhead\bottomrule\endfoot
allocator\_\allowbreak{}traits & The uniform adapter exposing allocator properties and operations.\\
pointer\_\allowbreak{}traits & A customization surface for pointer-like types.\\
unique\_\allowbreak{}ptr & An exclusive-ownership smart pointer.\\
shared\_\allowbreak{}ptr & A shared-ownership smart pointer with a control block.\\
weak\_\allowbreak{}ptr & A non-owning observer of a shared control block.\\
memory\_\allowbreak{}resource & A runtime-polymorphic allocation interface.\\
uses\_\allowbreak{}allocator & The protocol by which construction receives an allocator.\\
\end{longtable}
\CornellTakeaway{Safe memory management requires aligning storage acquisition, object lifetime, ownership, allocator propagation, and reclamation.}
\end{document}
